\documentclass[12pt]{article}
\usepackage[margin=1in]{geometry}
\usepackage{amsmath, amssymb, amsthm}
\usepackage{setspace}
\usepackage{parskip}
\usepackage{booktabs}
\usepackage{array}
\usepackage{hyperref}
\usepackage{cite}

\doublespacing

\theoremstyle{plain}
\newtheorem{theorem}{Theorem}
\theoremstyle{definition}
\newtheorem{definition}{Definition}
\newtheorem{example}{Example}

\title{\textbf{Vector Spaces, Linear Transformations, and LU Factorization\\ in Infectious Disease Modeling}}
\author{Charlene Miciano - EN609}
\date{June 23, 2026}

\begin{document}

\maketitle
\thispagestyle{empty}

\newpage

% -------------------------------------------------------
\section*{1. \quad What Is Compartmental Disease Modeling and Why Does It Matter?}
% -------------------------------------------------------

Infectious disease modeling is one of the most practically significant subfields of applied
mathematics. When a new pathogen emerges, public health authorities must answer a
deceptively simple question: \emph{will this disease spread through the population, or will
it die out?} Epidemiologists could not answer this question without the mathematical
foundations that lie in the linear algebra concepts we have studied in Modules 1--5:
\textbf{vector spaces}, \textbf{linear transformations}, \textbf{matrix inverses solved
via LU factorization}, and \textbf{elementary row operations}. With my background in
public health modeling, I have decided to dedicate this application paper to ``pull back
the curtain'' and truly understand the origins of one of the most used methods in
infectious disease modeling that I personally use in my own work in Tuberculosis modeling.

The dominant modeling framework is the \emph{compartmental model}, which divides a
population into mutually exclusive subgroups (``compartments'') and tracks the flow of
individuals between them over time. The canonical example is the SIR model, which
partitions every person into one of three states:

\begin{itemize}
  \item $S$ --- \textbf{Susceptible}: individuals who have not yet been infected and can
  acquire the disease.
  \item $I$ --- \textbf{Infectious}: individuals who are currently infected and can
  transmit the pathogen.
  \item $R$ --- \textbf{Recovered} (or Removed): individuals who have cleared the
  infection and are immune.
\end{itemize}

More realistic diseases require more compartments. Tuberculosis, for instance, features a
long latency period where a person harbors the bacterium without being infectious; this
necessitates an Exposed ($E$) state, giving the \textbf{SEIR} model. A TB cascade-of-care
model might further subdivide populations by treatment status, loss to follow-up (LTFU),
and drug resistance, producing state spaces with ten or more compartments.

The central question in each of these models is captured by a single number: the
\textbf{basic reproduction number} $R_0$. Informally, $R_0$ is the average number of
new infections caused by a single infected individual introduced into a fully susceptible
population \cite{diekmann2010}. The threshold criterion is as follows:
\[
  \text{if } R_0 < 1, \text{ the disease dies out}; \qquad
  \text{if } R_0 > 1, \text{ the disease spreads.}
\]

The matrix theory concepts that underly the computing of $R_0$ in structured populations is the
\textbf{next-generation matrix} (NGM), introduced formally by Diekmann, Heesterbeek,
and Metz in 1990 \cite{diekmann1990} and systematized for compartmental ODE models by
van den Driessche and Watmough in 2002 \cite{vddriessche2002}. Constructing the NGM
requires us to view the state space of infected individuals as a \textbf{vector space},
decompose the disease dynamics into two \textbf{linear transformations} ($F$ and $V$),
and invert $V$ using \textbf{LU factorization}. In staged disease models, this inversion
reduces to simply forward substitution on a lower-triangular system. The product
$K = FV^{-1}$ is the NGM, a linear transformation whose entries encode the expected
number of new infections each compartment generates.

New terms introduced in this paper: \textbf{compartmental model},
\textbf{disease-free equilibrium (DFE)}, \textbf{next-generation matrix (NGM)},
\textbf{basic reproduction number $R_0$}, \textbf{transmission matrix $F$},
\textbf{transition matrix $V$}, \textbf{M-matrix} (a square matrix with non-positive
off-diagonal entries whose determinant is strictly positive, guaranteeing invertibility),
and \textbf{Z sign pattern} (a matrix whose off-diagonal entries are all non-positive).

% -------------------------------------------------------
\section*{2. \quad The Mathematical Framework}
% -------------------------------------------------------

\subsection*{2.1 \quad The State Space as a Vector Space}

Let a disease model have $n$ total compartments. We represent the population's state at
time $t$ as a vector $\mathbf{x}(t) = (x_1, x_2, \ldots, x_n)^T \in \mathbb{R}^n_{\geq 0}$.
The first $m$ components correspond to \emph{infected} compartments (e.g., $E$ and $I$
in an SEIR model); the remaining $n - m$ components are \emph{uninfected} (e.g., $S$
and $R$). This infected subspace $\mathbb{R}^m \subset \mathbb{R}^n$ is the vector space
on which the NGM acts as a linear transformation. The fact that $\mathbf{x}(t)$ lives in
$\mathbb{R}^n_{\geq 0}$ reflects the real-world biological constraint (no negative people)
that imposes important structural properties on $F$ and $V$.

The dynamics of the full system are governed by a system of ODEs:
\[
  \dot{x}_i = \mathcal{F}_i(\mathbf{x}) - \mathcal{V}_i(\mathbf{x}), \quad i = 1, \ldots, n,
\]
where $\mathcal{F}_i(\mathbf{x})$ is the rate at which \emph{new infections} enter
compartment $i$, and $\mathcal{V}_i(\mathbf{x})$ captures all other flows: progression
through disease stages, recovery, death, and movement in and out via treatment
\cite{vddriessche2002}.

\subsection*{2.2 \quad Linearization at the Disease-Free Equilibrium}

The \textbf{disease-free equilibrium (DFE)}, denoted $\mathbf{x}^0$, is the point where
all infected compartments are empty. To study whether a small introduction of infectious
individuals will grow or decay, we linearize the infected subsystem at $\mathbf{x}^0$ by
taking partial derivatives. The resulting Jacobian of the infected subsystem decomposes as:
\[
  \left.\frac{\partial \mathbf{f}}{\partial \mathbf{x}}\right|_{\mathbf{x}^0,\,\text{infected}} = F - V,
\]
where $F = \left[\frac{\partial \mathcal{F}_i}{\partial x_j}(\mathbf{x}^0)\right]$ and
$V = \left[\frac{\partial \mathcal{V}_i}{\partial x_j}(\mathbf{x}^0)\right]$ are both
$m \times m$ matrices \cite{vddriessche2002}. Their key structural properties are:

\begin{itemize}
  \item $F$ is \textbf{non-negative}: every entry is a rate of new infection production,
  which must be $\geq 0$ by biology.
  \item $V$ is an \textbf{M-matrix}: it has non-positive off-diagonal entries (inflows from
  other infected classes) and positive diagonal entries (total outflows from each
  compartment). The conditions of the model force $\det(V) > 0$, so $V$ is
  \textbf{invertible}, and moreover $V^{-1} \geq 0$ entry-wise \cite{vddriessche2002}.
\end{itemize}

These two matrices are \textbf{linear transformations} on $\mathbb{R}^m$. The
transformation $F$ maps a current distribution of infected individuals to the rate of new
infections it generates; $V$ maps that same distribution to the rate at which individuals
transition out of infected states. Separating these two processes is the conceptual core of
the NGM method \cite{diekmann2010}.

\subsection*{2.3 \quad $V^{-1}$ as a Linear Transformation: Time-in-Compartment}

Before computing $K = FV^{-1}$, we should understand what $V^{-1}$ means as a linear
transformation, because its epidemiological interpretation is what makes the NGM readable.

With reinfection turned off, the infected subsystem evolves as $\dot{\mathbf{w}} = -V\mathbf{w}$,
whose solution is $\mathbf{w}(t) = e^{-Vt}\mathbf{w}(0)$. Integrating from $t = 0$ to
$\infty$ gives the total time each compartment is occupied:
\[
  \int_0^\infty e^{-Vt}\, dt = V^{-1}.
\]
This means the $(j,k)$ entry of $V^{-1}$ is the \textbf{expected total time} an individual
introduced into compartment $k$ spends in compartment $j$ over their entire infectious
career, barring reinfection \cite{vddriessche2002}. Combined with the fact that the $(i,j)$
entry of $F$ is the \textbf{rate} at which an individual in compartment $j$ generates new
infections in compartment $i$, the $(i,k)$ entry of $K = FV^{-1}$ is:
\begin{align*}
  K_{ik} = \,&\text{expected number of new infections in compartment } i \\
             &\text{generated by one individual seeded in compartment } k.
\end{align*}
In other words, $K = FV^{-1}$ is a linear transformation that maps a vector of initial
infection seeds to the vector of infections in the next generation \cite{diekmann1990}.

% -------------------------------------------------------
\section*{3. \quad Illustrative Example: The TB Treatment Model}
% -------------------------------------------------------

To make the framework concrete, we work through the SEIR-with-treatment model from
van den Driessche and Watmough \cite{vddriessche2002}, which directly mirrors the TB
models I work with. The population is divided into four compartments: $S$ (susceptible),
$E$ (exposed/latent), $I$ (infectious), and $T$ (treated), with infected compartments $E$
and $I$ giving $m = 2$. The parameters are: $\beta_1$ (transmission rate), $\nu$
(progression rate from $E$ to $I$), $r_1$ and $r_2$ (treatment rates for $E$ and $I$
respectively), $p = 1 - q$ (treatment failure probability), and $d$ (natural death rate).

Evaluating the partial derivatives at the DFE (with $S_0 = 1$) gives:
\[
  F = \begin{pmatrix} 0 & \beta_1 \\ 0 & 0 \end{pmatrix}, \qquad
  V = \begin{pmatrix} d + \nu + r_1 & -pr_2 \\ -\nu & d + r_2 \end{pmatrix}.
\]

The zero in the $(2,1)$ position of $F$ reflects that newly infected individuals enter $E$
only: no one is placed directly in $I$ upon infection. This means the only
\textbf{state-at-infection} is $E$ \cite{diekmann2010}, a structural observation that falls
directly out of examining which rows of $F$ are nonzero.

\textbf{Step 1: Compute $V^{-1}$ via row reduction.}

We augment $V$ with the identity and apply elementary row operations:
\[
  \left[\begin{array}{cc|cc}
    d+\nu+r_1 & -pr_2 & 1 & 0 \\
    -\nu       & d+r_2 & 0 & 1
  \end{array}\right].
\]
Let $\delta = \det(V) = (d+\nu+r_1)(d+r_2) - \nu p r_2$. Since $V$ is an M-matrix,
$\delta > 0$ \cite{vddriessche2002}. Row-reducing to the identity on the left yields:
\[
  V^{-1} = \frac{1}{\delta}
  \begin{pmatrix} d + r_2 & pr_2 \\ \nu & d + \nu + r_1 \end{pmatrix}.
\]
Every entry of $V^{-1}$ is non-negative (since all parameters are positive and
$\delta > 0$), confirming the time-in-compartment interpretation.

\textbf{Step 2: Form the NGM $K = FV^{-1}$.}

\[
  K = \begin{pmatrix} 0 & \beta_1 \\ 0 & 0 \end{pmatrix}
      \cdot
      \frac{1}{\delta}\begin{pmatrix} d + r_2 & pr_2 \\ \nu & d + \nu + r_1 \end{pmatrix}
    = \frac{1}{\delta}
      \begin{pmatrix} \beta_1\nu & \beta_1(d+\nu+r_1) \\ 0 & 0 \end{pmatrix}.
\]

\textbf{Step 3: Read $R_0$ from the structure of $K$.}

The second row of $K$ is zero, reflecting the fact that $I$ is not a state-at-infection:
the $I$ compartment produces infections, but no new infections are seeded directly into
$I$. Row-reducing $K$ (pivoting on the $(1,1)$ entry, which is the only nonzero pivot)
shows that $K$ has rank 1. When $K$ has rank 1, the NGM generates exactly one
independent infection pathway, and $R_0$ is the unique nonzero entry in the row-echelon
form of $K$:
\[
  \boxed{R_0 = \frac{\beta_1\nu}{\delta} = \frac{\beta_1\nu}{(d+\nu+r_1)(d+r_2) - \nu p r_2}.}
\]

\textbf{Interpretation.} The numerator $\beta_1\nu$ is the product of the transmission
rate and the rate at which latent individuals progress to infectiousness. The denominator
$\delta$ encodes the total competition between disease progression and all forms of outflow
(death, treatment, treatment failure cycling). As treatment improves through larger $r_1$,
$r_2$, or $q$, $\delta$ increases, $R_0$ decreases, and the linear transformation $K$
shrinks the next generation. In other words, better treatment reduces epidemic potential.

% -------------------------------------------------------
\section*{4. \quad Advanced Example: LU Factorization and the Staged Progression Model}
% -------------------------------------------------------

\subsection*{4.1 \quad Structure of the Staged Progression Model}

The TB treatment example above used a $2 \times 2$ matrix $V$ that we inverted by row
reduction. In a \textbf{staged progression model}, infected individuals move sequentially
through $m$ disease stages with changing infectivity. This framework applies directly to
HIV, TB drug-resistance progression, and multi-stage cascade-of-care models. In this
setting, the matrix $V$ takes on a lower-bidiagonal structure that is solved by forward
substitution, which is the solve step of \textbf{LU factorization} with $L = V$ and $U = I$.

The model has a single susceptible compartment $S$ and $m$ infectious stages
$I_1, I_2, \ldots, I_m$, with progression rate $\nu_i$ out of stage $i$ and
disease-induced death rate $d_i$ at stage $i$ \cite{vddriessche2002}. The infected
subsystem linearized at the DFE gives:
\[
  F = \begin{pmatrix}
    \beta_1 & \beta_2 & \cdots & \beta_{m-1} & 0 \\
    0       & 0       & \cdots & 0           & 0 \\
    \vdots  &         &        &             & \vdots \\
    0       & 0       & \cdots & 0           & 0
  \end{pmatrix},
  \qquad
  V = \begin{pmatrix}
    \nu_1+d_1 & 0         & 0         & \cdots & 0   \\
    -\nu_1    & \nu_2+d_2 & 0         & \cdots & 0   \\
    0         & -\nu_2    & \nu_3+d_3 & \cdots & 0   \\
    \vdots    &           & \ddots    & \ddots & \vdots \\
    0         & 0         & \cdots    & -\nu_{m-1} & d_m
  \end{pmatrix},
\]
where the final stage $I_m$ has zero infectivity ($\beta_m = 0$) due to morbidity, and
$\nu_m := 0$ is set for notational convenience \cite{vddriessche2002}. Because $V$ is
already lower triangular, it is its own $L$ factor in an LU decomposition with $U = I$.

\subsection*{4.2 \quad Lemma: $F$ is Non-Negative and $V$ is an Invertible M-Matrix}

The following result from van den Driessche and Watmough \cite{vddriessche2002} (Lemma 1)
establishes the two structural properties that make the NGM well-defined. All of the
epidemiological assumptions cited below are stated formally in that paper; here we describe
them in plain terms.

\begin{theorem}[Lemma 1, van den Driessche \& Watmough \cite{vddriessche2002}]
Let $\mathbf{x}^0$ be a DFE of the compartmental model. Under the biological assumptions
of the model (non-negativity of population counts, no immigration of infectives, and
stability of the DFE in the absence of infection), the matrices $F$ and $V$ defined by the
linearization at $\mathbf{x}^0$ satisfy: (i) $F \geq 0$ entry-wise, (ii) $V$ is a nonsingular
M-matrix, and (iii) $V^{-1} \geq 0$ entry-wise.
\end{theorem}

\begin{proof}[Proof summary, following van den Driessche \& Watmough \cite{vddriessche2002}]
We summarize the three conclusions in turn.

\textbf{(i) $F \geq 0$.} The function $\mathcal{F}_i(\mathbf{x})$ counts new infections
entering compartment $i$, so it is non-negative for all $\mathbf{x} \geq 0$ by biology. At
the DFE, the biological assumptions also require $\mathcal{F}_i(\mathbf{x}^0) = 0$ for
infected compartments (since there is no disease present to generate new infections). The
entry $F_{ij} = \frac{\partial \mathcal{F}_i}{\partial x_j}(\mathbf{x}^0)$ is a one-sided
limit of the form $\lim_{h \to 0^+} [\mathcal{F}_i(\mathbf{x}^0 + h\mathbf{e}_j) -
\mathcal{F}_i(\mathbf{x}^0)]/h$. Because $\mathcal{F}_i \geq 0$ and
$\mathcal{F}_i(\mathbf{x}^0) = 0$, this limit is $\geq 0$. Thus $F \geq 0$.

\textbf{(ii) $V$ is a nonsingular M-matrix.} We show $V$ has the Z sign pattern and
positive diagonal, which together with the stability assumption imply $V$ is a nonsingular
M-matrix. For the off-diagonal entries: if compartment $i$ is empty ($x_i = 0$), no
individuals can leave it, so $\mathcal{V}^-_i(\mathbf{x}) = 0$. The entry $V_{ij}$ for
$i \neq j$ is the partial derivative of net outflow from compartment $i$ with respect to
$x_j$. Because $x_i = 0$ forces $\mathcal{V}^-_i(\mathbf{x}^0 + h\mathbf{e}_j) \leq 0$,
we get $V_{ij} \leq 0$ for $i \neq j$, giving the Z sign pattern. The biological assumption
that the DFE is stable when $F = 0$ forces all eigenvalues of $-V$ to have negative real
parts, which is equivalent to $V$ having all eigenvalues with positive real parts. A matrix
with the Z sign pattern and positive real-part eigenvalues is by definition a nonsingular
M-matrix, so $\det(V) > 0$ and $V$ is invertible.

\textbf{(iii) $V^{-1} \geq 0$.} A standard result in the theory of M-matrices states that
if $V$ is a nonsingular M-matrix, then $V^{-1} \geq 0$ entry-wise \cite{vddriessche2002}.
Intuitively, this is consistent with the time-in-compartment interpretation: expected times
are non-negative.

This lemma is the structural foundation of the entire NGM framework. Once we know
$F \geq 0$, $V^{-1} \geq 0$, and $\det(V) > 0$, the product $K = FV^{-1}$ is
well-defined and non-negative, and its entries carry the infection-generation interpretation
described in Section 2.3.
\end{proof}

\subsection*{4.3 \quad Reading $R_0$ from $K = FV^{-1}$}

Since only the first row of $F$ is nonzero, $K = FV^{-1}$ also has only a nonzero first
row, and $R_0$ is the sum of the first-row entries of $K$. Using the closed-form entries
of $V^{-1}$ that follow from forward substitution on the lower-triangular $V$
\cite{vddriessche2002} (equation 12 therein):
\[
  R_0 = \sum_{j=1}^{m-1} \beta_j \cdot (V^{-1})_{1j}
      = \frac{\beta_1}{\nu_1+d_1}
        + \frac{\beta_2\nu_1}{(\nu_1+d_1)(\nu_2+d_2)}
        + \cdots
        + \frac{\beta_{m-1}\prod_{k=1}^{m-2}\nu_k}{\prod_{k=1}^{m-1}(\nu_k+d_k)}.
\]
Each term is transparent: the $j$-th term is the product of $\beta_j$ (infectivity in stage
$j$) and the fraction of individuals originally seeded in stage 1 who survive to reach stage
$j$, multiplied by the time they spend there. LU factorization via forward substitution on a
lower-triangular $V$ delivers this interpretation directly.

In a TB cascade-of-care model, adding more stages (diagnosed, on-treatment, LTFU,
re-treated) extends the size of $V$ but preserves its lower-bidiagonal structure. Each new
stage appends one row and one column, and forward substitution remains the efficient path
to $V^{-1}$. The NGM makes each intervention lever visible as a distinct entry in a linear
transformation, giving both modelers and policymakers a precise handle on which parameters
most influence $R_0$. After learning about and using compartmental models in public health
research for the past 3 years, it is incredible to finally understand how the methods were
developed and the matrix theory concepts that made it possible.

\newpage
\begin{thebibliography}{9}

\bibitem{diekmann1990}
O.\ Diekmann, J.\ A.\ P.\ Heesterbeek, and J.\ A.\ J.\ Metz,
``On the definition and the computation of the basic reproduction ratio $R_0$ in
models for infectious diseases in heterogeneous populations,''
\textit{Journal of Mathematical Biology}, vol.\ 28, pp.\ 365--382, 1990.

\bibitem{diekmann2010}
O.\ Diekmann, J.\ A.\ P.\ Heesterbeek, and M.\ G.\ Roberts,
``The construction of next-generation matrices for compartmental epidemic models,''
\textit{Journal of the Royal Society Interface}, vol.\ 7, no.\ 47, pp.\ 873--885, 2010.

\bibitem{vddriessche2002}
P.\ van den Driessche and J.\ Watmough,
``Reproduction numbers and sub-threshold endemic equilibria for compartmental models
of disease transmission,''
\textit{Mathematical Biosciences}, vol.\ 180, pp.\ 29--48, 2002.

\end{thebibliography}

\end{document}
